\Page{19}{\textenglish{LAST LOOK}}{ארבע החלטות שהשתנו}{55}{63}

\vspace{9mm}

{\fontsize{17.3}{26.2}\selectfont
סטודיו יכול לראות יותר גרסאות לפני שיבחר מה לתפור. קניינית יכולה לאתר סימן מוקדם לצבע עולה ולבדוק אותו בהזמנה קטנה. בשני המקרים, היכולת להוסיף אפשרויות מעניינת בעיקר כשהארגון מוכן להשאיר זמן וכסף לניסוי. אם כל צוות פונה לאותות שכבר נראים חזקים, אפילו תחזית מוצלחת עלולה לדחוף קולקציות שונות לעבר אותן בחירות.

בהפקה ובקנייה עומד למבחן סוג אחר של ודאות. אפשר לערוך דימוי שוב ושוב, אבל צריך להחליט מי רשאי להשתמש בפנים שמופיעות בו ומה ראוי לגלות לקהל. למעשה, האמון שהקונה נותנת בתמונה תלוי גם במה שהתמונה מבטיחה לה. הצצה לגזרה ומענה על השאלה איזו מידה להזמין מצריכים בדיקות שונות.

\par}

\vspace{5mm}
{\color{Silver}\hrule height .55pt}
\vspace{5mm}

{\fontsize{17.3}{26.2}\selectfont
מבחן טוב לטכנולוגיה הוא לברר איזו החלטה נעשתה קלה יותר לתיקון. אפשר לנסות להזמין כמות קטנה לפני שמרחיבים ייצור, או לבדוק דוגמה שנראתה מבטיחה במסך לפני שמאשרים אותה. כשמודדים הצלחה רק במספר ההדמיות או במהירות שבהן נוצרו, לא יודעים אם הניסוי שיפר את הקולקציה או את מה שיישאר ממנה במחסן.

יש מקום לאלגוריתם בבית אופנה שמסוגל להסביר מדוע קיבל המלצה, למדוד מה קרה בעקבותיה ולשמור גם תקציב לרעיון שהנתונים טרם למדו לזהות.
\par}
